\section{Related Works}
\subsection{Task Specification in Manipulation}

Specifying tasks for the wide range of manipulation problems spans symbolic, learning-based, outcome-driven, and object-centric interfaces. Classical approaches encode goals and constraints with symbolic formalisms such as PDDL and temporal logics, or optimize cost-augmented formulations \cite{garrett2021integrated,zhao2024survey,toussaint2015logic}. Learning-based systems specify tasks often through language and perception, mapping instructions to actions via language-conditioned visuomotor policies and vision–language–action models \cite{shridhar2022cliport,shridhar2023perceiver,zitkovich2023rt,black2024pi_0,kim2025openvla}. Outcome-based specification sets goals by example observations, e.g., image goals with goal-conditioned policies \cite{lynch2020learning,ebert2018visual,black2023zero,xie2018few,sharma2019third,mendonca2021discovering}, and some works incorporate force targets \cite{adeniji2025feel}. Object-centric alternatives rely on descriptors or keypoints to capture task-relevant structure~\cite{simeonov2022neural}. Recently, foundation models enable higher-level interfaces that compile intent into actionable specifications via code~\cite{liang2023code}, 3D value maps akin to potential fields~\cite{huang2023voxposer}, keypoint relations~\cite{huang2025rekep}, or affordance maps~\cite{tang2025uad}.

\subsection{2D/3D Flow in Robotics}

Dense motion fields—optical flow, point tracks, and 3D scene/object flow—provide an embodiment-agnostic, mid-level interface for manipulation \cite{xu2025flow,yuan2025general}. In scene-centric formulations, policies parameterize or condition on motion in 2D or 3D to decide actions, with point/keypoint interfaces unifying perception and action and with action-flow improving precision \cite{weng2021fabricflownet,goyal2022ifor,seita2023toolflownet,chen2024g3flow,wang2025skil,haldar2025point,guo2025actionsink}. In object-centric formulations, desired object motion is specified independently of embodiment and then converted into actions via policy inference, planning, and optimization \cite{eisner2022flowbot3d,gao2024flip,guo2025flowdreamer,zhi2025flowaction,he2025manitrend,yin2025object}. When actions are absent, retargeting with predicted tracks and dense correspondences offers a practical bridge across embodiments \cite{bharadhwaj2024track2act,ko2023learning,chen2025ecflow}. Advances in perception and tracking such as RGB-D motion-based segmentation, 3D scene flow in point clouds, zero-shot monocular and ToF-based scene flow, rigid-motion learning, refractive flow for transparent objects, deformable and articulated reconstruction, and structured scene representations \cite{shao2018motionseg,liu2018flownet3d,liang2025zeroshotmsf,nvidia2025zeromsf,sander2025tofsceneflow,byravan2016se3nets,tang2023rftrans,duisterhof2024deformgs,kerr2024rsrd,jiang2024roboexp,shorinwa2024splatmover} make the interface reliable. Flow-derived supervision further supports learning, and video generation supplies plausible visual rollouts for planning and imitation \cite{guzey2025bridging,yu2025genflowrl,bharadhwaj2024gen2act,patel2025rigvid}. Our approach follows the object-centric path by reconstructing 3D object flow from language-conditioned generations and tracking it under embodiment constraints, complementing other flow-conditioned policy representations~\cite{weng2021fabricflownet,seita2023toolflownet,eisner2022flowbot3d,gao2024flip,guo2025flowdreamer}.

\subsection{Video Models for Robotics}
Recent work increasingly integrates video models across robotic tasks in various ways \cite{yang2024video,mccarthy2025towards}. They can serve as auxiliary training objectives \cite{wu2023unleashing,seo2022reinforcement,wu2023pre,yang2024spatiotemporal,hu2024video,li2025unified}, as reward models \cite{huang2024diffusion,escontrela2023video,chen2021learning}, as policies \cite{ajay2023compositional,du2023learning}, or as a simulator for the environments \cite{valevski2024diffusion,bruce2024genie}. 
Notably, predictive modeling in robotics can leverage video frame prediction as a form of world model. By simulating future visual observations, these models can enable visual planning and manipulation by anticipating how the environment will evolve. For example, a video generative model can simulate long-horizon task outcomes, effectively acting as a visual planner \cite{luo2024grounding}. Video models in this way can also serve as 
as dynamics models \cite{du2023learning,yang2023learning,zhou2024robodreamer,rybkin2018learning,mendonca2023structured,che2024gamegen}.
Another notable direction is that video generation can directly provide new training data for robot learning. Several recent works devise frameworks to imagine new trajectories in the form of videos and use them to train or finetune policies, particularly for imitation learning \cite{jang2025dreamgen,liang2024dreamitate}.
